\documentclass[10pt,a4paper]{article}
\usepackage{amsmath}
\usepackage{amssymb}
\usepackage[utf8]{inputenc}
\usepackage{geometry}
\geometry{margin=1in}

\begin{document}
	
	\section*{Simple Model Comparison Example}
	
	\subsection*{Data}
	Five coin tosses resulted in 4 heads and 1 tail.
	
	\subsection*{Models}
	\begin{itemize}
		\item \(M_1\): The coin’s probability of heads is \(p = 0.5\).
		\item \(M_2\): The coin’s probability of heads is \(p = 0.8\).
	\end{itemize}
	
	Assume equal prior probabilities:
	\[
	P(M_1) = P(M_2) = 0.5.
	\]
	
	\subsection*{Likelihoods}
	\[
	P(\text{4 heads,1 tail} \mid M_1)
	= \binom{5}{4}\,(0.5)^4\,(0.5)^1
	= 5 \times 0.5^5 = 0.15625.
	\]
	\[
	P(\text{4 heads,1 tail} \mid M_2)
	= \binom{5}{4}\,(0.8)^4\,(0.2)^1
	= 5 \times 0.8^4 \times 0.2 = 0.4096.
	\]
	
	\subsection*{Posterior Model Probabilities}
	\[
	P(M_1 \mid D)
	= \frac{0.15625 \times 0.5}
	{0.15625 \times 0.5 + 0.4096 \times 0.5}
	\approx 0.276.
	\]
	\[
	P(M_2 \mid D)
	= \frac{0.4096 \times 0.5}
	{0.15625 \times 0.5 + 0.4096 \times 0.5}
	\approx 0.724.
	\]
	
	\subsection*{Interpretation}
	Given the data and equal prior model probabilities, model \(M_2\) (with \(p = 0.8\)) is about 72.4\% likely, whereas model \(M_1\) (with \(p = 0.5\)) is about 27.6\% likely.  
	Thus we conclude: *“It is more plausible that the coin has \(p=0.8\) rather than \(p=0.5\) under our assumptions and the observed data.”*
	
\end{document}
